\chapter[Desenvolvimento do Backend]{Desenvolvimento do Backend} \label{capitulo8}

O backend do sistema foi desenvolvido utilizando Node.js como ambiente de execução e Express.js como framework para construção da API REST. A camada de persistência utiliza o Prisma ORM, com banco de dados MySQL, enquanto a autenticação dos usuários é realizada por meio de JSON Web Tokens (JWT), utilizando tokens de acesso e de renovação. Para a validação dos dados recebidos pela API foi utilizado o Zod, e a documentação dos recursos disponibilizados é realizada por meio do padrão Swagger/OpenAPI.

A aplicação foi organizada em módulos de negócio, buscando separar as responsabilidades relacionadas às diferentes funcionalidades do sistema. A estrutura do backend contém módulos específicos para autenticação, usuários, aprendentes, avaliações, aplicações, relatórios e dashboard. Além desses módulos, foram definidos diretórios específicos para configurações, middlewares e utilitários, conforme apresentado na estrutura do projeto.

Essa organização permite concentrar em cada módulo as operações relacionadas a um determinado domínio da aplicação. O módulo de autenticação, por exemplo, concentra as funcionalidades relacionadas ao acesso ao sistema, enquanto os módulos de aprendentes, avaliações e aplicações estão relacionados diretamente ao processo de avaliação psicopedagógica. O módulo de relatórios é responsável pelas operações relacionadas às devolutivas produzidas a partir das aplicações, e o módulo de dashboard disponibiliza informações destinadas à visualização geral dos dados do sistema.

\section{Acesso aos dados}

A comunicação entre o frontend e o backend é realizada por meio de uma API REST, disponibilizada pelo Express.js. Os recursos são organizados em diferentes rotas, de acordo com os módulos funcionais do sistema.

Entre os principais recursos disponibilizados estão as operações de autenticação, gerenciamento de usuários, gerenciamento de aprendentes, gerenciamento de avaliações, aplicações, relatórios e dashboard. O endpoint de autenticação disponibiliza a operação \textit{POST /api/auth/login}, enquanto os demais módulos possuem seus respectivos prefixos de rota.

O gerenciamento de usuários é disponibilizado por meio das operações \textit{GET}, \textit{POST}, \textit{PUT} e \textit{DELETE} em \textit{/api/users}. De maneira semelhante, o módulo de aprendentes disponibiliza operações de consulta, cadastro, alteração e exclusão por meio de \textit{/api/aprendentes}, enquanto as avaliações são gerenciadas por meio de \textit{/api/avaliacoes}. Para aplicações, são disponibilizadas operações de consulta e criação em \textit{/api/aplicacoes}, e os relatórios possuem operações correspondentes em \textit{/api/relatorios}. O sistema também disponibiliza o endpoint \textit{GET /api/dashboard} para acesso às informações utilizadas no painel da aplicação.

A utilização de uma API com recursos separados por domínio permite que o frontend consuma as funcionalidades do sistema de maneira padronizada, mantendo a lógica de processamento e acesso aos dados concentrada no backend.

\section{Organização dos módulos e middlewares}

A estrutura interna da aplicação está dividida entre configurações, middlewares, módulos de negócio e utilitários. O diretório \textit{config} concentra configurações relacionadas ao ambiente, ao Prisma e ao Swagger. Já o diretório \textit{middlewares} contém mecanismos relacionados à autenticação, validação e tratamento de erros.

Os middlewares são utilizados como uma camada intermediária no processamento das requisições. Dessa forma, a autenticação pode ser aplicada antes do acesso a recursos protegidos, os dados recebidos podem ser submetidos às regras de validação e erros ocorridos durante o processamento podem ser tratados de maneira centralizada.

A validação das informações de entrada é realizada com o Zod. Essa abordagem permite definir as estruturas esperadas para os dados recebidos pela API e verificar se as informações fornecidas estão de acordo com essas estruturas antes de serem processadas pela aplicação.

\section{Autenticação e controle de acesso}

A autenticação foi implementada utilizando JWT, com utilização de tokens de acesso e de renovação (refresh token). Essa estratégia permite identificar o usuário autenticado durante o acesso aos recursos protegidos da API e manter o gerenciamento da sessão por meio dos diferentes tipos de token.

A autenticação está associada ao módulo \textit{auth} e ao middleware correspondente, responsável por participar do processo de proteção das rotas da aplicação. O backend também contempla diferentes perfis de usuário, definidos no modelo de dados como \textit{ADMIN}, \textit{PSICOPEDAGOGO} e \textit{VIEWER}, permitindo estabelecer diferentes níveis de acesso às funcionalidades do sistema. 

Além do controle por perfil, o sistema implementa restrição de visibilidade baseada em titularidade: aprendentes e avaliações são vinculados ao usuário que os criou, e as operações de leitura retornam somente os registros pertencentes ao usuário autenticado, exceto para o perfil \textit{ADMIN}, que possui acesso irrestrito.

Além da autenticação, o backend possui mecanismos de validação das requisições e tratamento de erros. A combinação desses mecanismos contribui para que as operações disponibilizadas pela API sejam executadas somente após o processamento das informações necessárias para cada recurso.

\section{Gerenciamento de aprendentes e avaliações}

O módulo de aprendentes disponibiliza as operações necessárias para o cadastro e gerenciamento dos indivíduos acompanhados pelo profissional. Os dados são posteriormente utilizados nas aplicações dos instrumentos avaliativos, mantendo a associação entre o aprendente e os registros produzidos durante o processo de avaliação.

A visibilidade dos aprendentes é controlada pelo campo \textit{responsavelId}, preenchido automaticamente com o identificador do usuário autenticado no momento do cadastro. Nas operações de listagem e consulta individual, o módulo aplica um filtro que restringe os resultados ao usuário responsável. O perfil \textit{ADMIN} constitui exceção a essa regra, tendo acesso irrestrito a todos os registros. Essa lógica é implementada na camada de serviço, antes da consulta ao banco de dados, sem expor ao cliente a distinção entre registros próprios e de outros usuários.

O módulo de avaliações é responsável pelo gerenciamento dos instrumentos utilizados pelo sistema. A estrutura adotada permite representar avaliações compostas por diferentes questões e tipos de resposta. O versionamento desses instrumentos é mantido na camada de persistência, conforme apresentado na modelagem do banco de dados, permitindo relacionar diferentes versões de uma mesma avaliação.

De maneira análoga ao módulo de aprendentes, a visibilidade das avaliações é controlada pelo campo \textit{autorId}, atribuído automaticamente ao usuário que criou o instrumento. As consultas de listagem e de detalhe aplicam o mesmo filtro de restrição por responsável, garantindo que cada profissional acesse somente os instrumentos que criou. Ao criar uma nova versão de uma avaliação, o campo \textit{autorId} é propagado para o novo registro, mantendo a consistência da titularidade ao longo do histórico de versões.

As aplicações são tratadas separadamente pelo módulo \textit{aplicacoes}, permitindo associar uma avaliação a um aprendente e registrar o andamento desse processo. Essa separação entre avaliação e aplicação possibilita que um mesmo instrumento seja utilizado em diferentes processos avaliativos, mantendo cada aplicação como um registro independente.

\section{Relatórios e dashboard}

O módulo \textit{relatorios} concentra as operações relacionadas aos relatórios produzidos a partir das aplicações realizadas. Esses registros permanecem associados à aplicação correspondente, permitindo relacionar a devolutiva produzida pelo profissional ao processo avaliativo que lhe deu origem.

O backend também possui um módulo específico de \textit{dashboard}, disponibilizado por meio da rota \textit{GET /api/dashboard}. Esse módulo fornece os dados necessários para a apresentação das informações gerais do sistema na interface, evitando que o frontend precise realizar diretamente operações sobre o banco de dados.

Nesse contexto, o backend atua como intermediário entre os dados persistidos e a interface do sistema, sendo responsável pelo processamento das requisições e pela disponibilização das informações necessárias às funcionalidades implementadas.

\section{Documentação e testes}

A API possui documentação baseada em Swagger/OpenAPI, disponibilizada durante a execução do sistema no endereço \textit{/api-docs}. A documentação permite consultar os recursos disponibilizados pela API e auxilia na compreensão e utilização dos endpoints durante o desenvolvimento e os testes.

Para os testes automatizados do backend foram utilizadas as ferramentas Vitest e Supertest. O projeto disponibiliza comandos específicos para execução dos testes e para geração de informações de cobertura, permitindo verificar o comportamento das funcionalidades implementadas.

Também foram definidos procedimentos para geração do cliente do Prisma, execução das migrações e população inicial do banco de dados. Essas operações permitem preparar o ambiente de desenvolvimento de forma padronizada e manter a estrutura do banco sincronizada com o modelo utilizado pela aplicação.

A rastreabilidade das informações é mantida por meio da estrutura do banco de dados, sem a utilização de um mecanismo de log de auditoria separado. Todas as entidades relevantes possuem campos \textit{createdAt} e \textit{updatedAt}, que registram automaticamente os momentos de criação e atualização de cada registro. As exclusões são realizadas de forma lógica por meio do campo \textit{deletedAt}, preservando o histórico dos dados sem removê-los fisicamente. Adicionalmente, os vínculos de autoria presentes no modelo — \textit{responsavelId} em aprendentes, \textit{autorId} em avaliações e relatórios, e \textit{aplicadorId} em aplicações — permitem identificar o usuário associado a cada registro. Essa abordagem possibilita reconstituir a autoria e a temporalidade das operações diretamente a partir dos dados persistidos.

Dessa forma, o desenvolvimento do backend estabeleceu a camada responsável pela disponibilização da API REST, autenticação, validação das requisições, aplicação das operações relacionadas aos módulos do sistema, comunicação com o banco de dados e disponibilização dos dados utilizados pelo frontend. A organização modular adotada também fornece uma estrutura para a manutenção e evolução das funcionalidades relacionadas ao processo de avaliação psicopedagógica.